\section{Concluding Remarks} \label{section:concluding_remarks}

% ---------- %

This work demonstrates that while general-purpose MLPs like MACE can successfully preserve interface rankings in
selected systems, their predictive utility is highly system- and geometry-dependent.

Strong performance was observed in Ge, Sn, and C|X systems with good atomic registry, as well as in chemically uniform
Si|Ge alloy configurations.

In contrast, Si-rich and poorly coordinated interfaces exhibited systematic rank failures, underscoring the limits of
current models in strained or covalently dominated regimes.

% ---------- %

By enabling fast, large-scale screening of interfacial structures in favourable systems, this study lays the foundation
for accelerated discovery workflows.

However, widespread application across diverse material combinations will require targeted retraining, delta-learning,
or integrated MLP-DFT hybrid schemes.

Future work should continue to refine these approaches, using rank-preservation metrics as a primary benchmark for
practical utility in materials design.

% ---------- %

Despite the promising results in rank preservation for certain systems, this study has also highlighted substantial
challenges in building a truly generalisable machine-learned potential for interface prediction.

Limitations such as incomplete training coverage, extrapolation to novel configurations, and model underperformance in
strained or covalently bonded systems underscore the need for further refinement.

These risks are being actively addressed through targeted retraining, delta-learning corrections, and expanded benchmark
testing, as outlined in Chapter~\ref{chapter:next_steps}.

While strong performance was observed in Ge, Sn, and alloy systems, achieving consistent accuracy across a broader range
of interfaces will require careful design of training data and model architecture.

% ---------- %
